\subsection{Numerical track-light references for the front-page table}
\label{sec:cosmicnumbers}

For an order-of-magnitude conversion, use a minimum-ionizing track loss
$w_\mu=\SI{210.5}{MeV/m}$ (PDG liquid-argon table~\cite{pdg_argon}) and
$Y_{\gamma,\mu}=20{,}000\ \mathrm{photons/MeV}$. If the time-averaged
track-length density $\xi_\mu$ is spatially uniform, its photon production
 density and VD channel rate are
\begin{equation}
q_{\gamma,\mu}=\xi_\mu w_\mu Y_{\gamma,\mu},\qquad
R_{\mu,\mathrm{ch}}^{\mathrm{VD}}
=q_{\gamma,\mu}G_{\mathrm{XA}}\epsilon_{\mathrm{XA}}/2.
\label{eq:uniformmu}
\end{equation}
Here $[\xi_\mu]=\mathrm{m\ of\ track}/(\mathrm{m^3\,s})$.
This uses the same optical integrals as $^{39}$Ar; spatial uniformity,
constant track loss, and negligible shower light are additional assumptions.

\paragraph{ProtoDUNE-VD surface.}
The PDG quotes approximately one muon/cm$^2$/min through a horizontal
 detector and an approximate $\cos^2\theta$ angular distribution
\cite{pdg_cosmic}. With intensity defined per area normal to the track,
$F_h=\int I(\theta)\cos\theta\,d\Omega$ whereas
$\xi_\mu=\int I(\theta)\,d\Omega=(4/3)F_h$. Hence
$F_h=\SI{166.7}{m^{-2}.s^{-1}}$ and
$\xi_\mu=\SI{222.2}{m^{-2}.s^{-1}}$. This gives about
$1.2\times10^7$ PE/s/cathode channel and $6.1\times10^6$ PE/s/wall channel.
These are generic unshielded surface values; the CERN overburden, angular distribution,
and CRT-measured track population must replace this reference for site rates.

\paragraph{FD-VD, deep-underground track reference.}
The measured total muon flux at the SURF Davis Campus is
$(5.31\pm0.17)\times10^{-9}$ cm$^{-2}$s$^{-1}$~\cite{surf_muons}.
For an order-of-magnitude track benchmark take
$\xi_\mu\simeq\SI{5.31e-5}{m^{-2}.s^{-1}}$, giving
4--9 PE/s/channel. Equating this site's scalar flux to track-length density
is an approximation: a DUNE-hall angular/overburden calculation supplies the
actual conversion. High-energy radiative losses and showers are excluded;
these numbers describe the minimum-ionizing track term only.

\paragraph{Full ND-LAr and the 2$\times$2 underground prototype.}
For a total detector crossing rate $f_\mu$ and mean \emph{active-argon}
track length $\overline\ell_\mu$, the uniform wall-flux approximation gives
\begin{equation}
\boxed{R_{\mu,\mathrm{ch},t}^{\mathrm{ND}}
=\frac{f_\mu\overline\ell_\mu w_\mu Y_{\gamma,\mu}
h_t\epsilon_t}{N_{\mathrm{ch},t}},\quad t=A,L.}
\label{eq:NDmu}
\end{equation}
The ND CDR data-volume model uses an anticipated 100 Hz cosmic crossing
rate for full ND-LAr~\cite{nd_cdr}. Adopt that rate and an \emph{assumed}
3 m mean active path: ArCLight/LCM give 90/271 PE/s/channel and the
all-channel mean is 180. Scale these entries by
$(f_\mu/100\ \mathrm{Hz})(\overline\ell_\mu/3\ \mathrm m)$.
The 2$\times$2 operation paper estimates approximately 2 Hz in MINOS Hall
\cite{nd_operation}. With an \emph{assumed} 1 m mean active path, its
ArCLight/LCM values are 12/39 PE/s/channel and the mean is 26. Scale by
$(f_\mu/2\ \mathrm{Hz})(\overline\ell_\mu/1\ \mathrm m)$.
These are different underground sites and sizes; neither the prototype's
2 Hz nor ProtoDUNE's surface flux is transferred to full ND-LAr.
Beam neutrinos, beam-induced rock muons, and associated secondaries require
 a separate spill-dependent contribution and are absent from this comparison.

The reference light sums are therefore about $10^7$ PE/s/channel for
surface ProtoDUNE, $10^5$ for FD-VD, a few $10^2$ for full ND-LAr, and a
few $10^1$ for underground 2$\times$2. These totals include only the two
modeled scintillation sources; they do not predict electronic trigger rates.
